\documentclass{article}
\usepackage{amsmath, amssymb, amsthm, geometry}
\geometry{a4paper, margin=1in}

\title{ACM India Summer School on Symmetric Key Cryptography \\ Day 3: Implementation of Block Ciphers (AES-128 and DES)}
\author{Sourav Sharma}
\date{June 10, 2026}

\begin{document}

\maketitle

\section{Introduction}
Block ciphers are a cornerstone of symmetric key cryptography. In this lab, I implemented the complete encryption and decryption pipelines for two major block ciphers: the Data Encryption Standard (DES) and the Advanced Encryption Standard (AES-128). Both implementations were written in clean, comment-free C++ (using a Competitive Programming style) and verified using standard test vectors.

\section{Data Encryption Standard (DES)}

DES is a symmetric block cipher that operates on 64-bit blocks of data using a 56-bit key. It relies on a Feistel network structure consisting of 16 rounds of substitution and permutation.

\subsection{Key Schedule}
The 64-bit master key $K$ contains 8 parity bits, which are discarded. The remaining 56 bits are permuted and split into two 28-bit registers, $C_0$ and $D_0$, using the Permuted Choice 1 (PC-1) table.
In each round $i \in \{1, \dots, 16\}$, the registers are left circularly shifted according to a predefined shift schedule (either 1 or 2 positions):
\[ C_i = \text{LS}_i(C_{i-1}), \quad D_i = \text{LS}_i(D_{i-1}) \]
The concatenated 56-bit state $C_i \parallel D_i$ is then compressed to a 48-bit round subkey $K_i$ using the Permuted Choice 2 (PC-2) table.

\subsection{Feistel Round Function}
The Feistel function $f(R, K_i)$ takes the 32-bit right half of the state $R$ and the 48-bit subkey $K_i$:
\begin{enumerate}
    \item \textbf{Expansion ($E$)}: The 32-bit block is expanded to 48 bits using the expansion table $E$, duplicating certain bits.
    \item \textbf{Key Addition}: The expanded block is XORed with the 48-bit subkey: $X = E(R) \oplus K_i$.
    \item \textbf{Substitution ($S$-boxes)}: The 48-bit result is split into eight 6-bit groups. For each group $j \in \{1, \dots, 8\}$, the outer bits determine the row, and the inner bits determine the column of the corresponding $S$-box $S_j$. Each $S$-box maps 6 bits to 4 bits.
    \item \textbf{Permutation ($P$)}: The concatenated 32-bit output of the S-boxes is permuted using the table $P$.
\end{enumerate}

\subsection{Encryption and Decryption Structure}
The input block is first permuted via the Initial Permutation (IP). The state is split into $L_0$ and $R_0$. For $i = 1$ to 16, the round updates are:
\[ L_i = R_{i-1}, \quad R_i = L_{i-1} \oplus f(R_{i-1}, K_i) \]
After the 16th round, the final state is swapped ($R_{16} \parallel L_{16}$) and permuted using the Final Permutation (FP), which is the inverse of IP. Decryption uses the exact same structure but applies the subkeys in reverse order ($K_{16}, K_{15}, \dots, K_1$).

\subsection{Verification}
I compiled and ran \texttt{DES\_Solution.cpp} against \texttt{DES\_TestVectors.csv}. The implementation successfully passed all \textbf{100/100 test vectors} (containing weak keys, semi-weak keys, and random inputs).

\section{Advanced Encryption Standard (AES-128)}

AES-128 is a SP-network block cipher that processes 128-bit blocks using a 128-bit key. The block state is formatted as a $4 \times 4$ byte grid (column-major order).

\subsection{Key Schedule}
The 128-bit master key is expanded into 44 words ($w_0, \dots, w_{43}$) of 32 bits each, providing 11 round keys (each of 128 bits). For indices not divisible by 4, $w_i = w_{i-4} \oplus w_{i-1}$. For $i$ divisible by 4, a non-linear transformation is applied to $w_{i-1}$ before XORing:
\[ \text{SubWord}(\text{RotWord}(w_{i-1})) \oplus \text{Rcon}[i/4] \]
where $\text{RotWord}$ performs a 1-byte left circular shift, $\text{SubWord}$ applies the AES S-box to each of the 4 bytes, and $\text{Rcon}$ contains the round constants in $\text{GF}(2^8)$.

\subsection{Round Operations}
AES-128 performs 10 rounds of four mathematical operations:
\begin{enumerate}
    \item \textbf{SubBytes}: A non-linear substitution step where each byte is replaced with its sub-byte using the AES S-box (constructed from the multiplicative inverse over $\text{GF}(2^8)$ followed by an affine transformation).
    \item \textbf{ShiftRows}: A transposition step where the rows of the state grid are cyclically shifted left by offset $r \in \{0, 1, 2, 3\}$.
    \item \textbf{MixColumns}: A linear transformation step where each column of the state is multiplied by a fixed polynomial over $\text{GF}(2^8)$ modulo $x^4 + 1$. The multiplication by 2 is implemented using a left shift and a conditional XOR with \texttt{0x1b} (the reduction polynomial for $\text{GF}(2^8)$: $x^8 + x^4 + x^3 + x + 1$).
    \item \textbf{AddRoundKey}: The state bytes are XORed with the corresponding round key bytes.
\end{enumerate}
In the final round (Round 10), the \text{MixColumns} step is omitted. Decryption is performed by applying the inverse operations (\text{InvSubBytes}, \text{InvShiftRows}, \text{InvMixColumns}) with the round keys in reverse order.

\subsection{Verification}
I compiled and tested \texttt{AES\_Solution.cpp} against \texttt{AES\_TestVectors.csv}. The implementation successfully passed all \textbf{100/100 test vectors}.

\end{document}
